\chapter{Теоретические сведения}
\label{ch:theory}

\section{Семантическая разметка HTML5}

\textbf{HTML5} вводит набор \textit{семантических} элементов, которые
описывают не оформление, а смысл блока: \texttt{<header>} (шапка),
\texttt{<main>} (основное содержимое), \texttt{<aside>} (вспомогательная
информация), \texttt{<section>} (раздел), \texttt{<article>} (самостоятельная
единица контента), \texttt{<nav>}, \texttt{<footer>}. Семантика улучшает
доступность (скринридеры понимают структуру), индексацию поисковыми
системами и читаемость кода.

Корректность разметки проверяется \textbf{валидатором W3C}
(\url{https://validator.w3.org}): он выявляет незакрытые теги, недопустимое
вложение элементов, отсутствие обязательных атрибутов (например,
\texttt{alt} у изображения). Для резюме семантическая структура естественна:
\texttt{<article>} --- всё резюме, \texttt{<aside>} --- боковая колонка с
контактами и навыками, \texttt{<section>} --- смысловые разделы (опыт,
образование), \texttt{<h1>}--\texttt{<h3>} --- иерархия заголовков.

\section{Редактируемый контент: атрибут \texttt{contenteditable}}

Глобальный атрибут \texttt{contenteditable="true"} превращает любой элемент в
редактируемое поле: пользователь кликает по тексту и правит его прямо на
странице, без отдельных \texttt{<input>} или \texttt{<textarea>}. Изменённый
текст остаётся \textit{внутри самого элемента} --- это удобно для
интерфейсов-конструкторов, где форма и результат совпадают.

\begin{lstlisting}[caption={Редактируемый элемент},captionpos=b]
<h1 contenteditable="true" data-field="name">Имя Фамилия</h1>
\end{lstlisting}

При работе с \texttt{contenteditable} учитываются особенности:
\begin{itemize}
  \item нажатие \texttt{Enter} по умолчанию вставляет блочный элемент
        (\texttt{<div>} или \texttt{<br>}) --- в однострочных полях это
        поведение перехватывается и отменяется;
  \item вставка из буфера (\texttt{paste}) может принести чужую разметку и
        стили --- содержимое буфера очищается до простого текста;
  \item опустевшее поле сохраняет служебный \texttt{<br>}, мешающий показу
        подсказки-плейсхолдера через псевдокласс \texttt{:empty}.
\end{itemize}

\section{Клиентское хранилище \texttt{localStorage}}

\textbf{Web Storage API} предоставляет объект \texttt{localStorage} ---
постоянное хранилище пар <<ключ--значение>> в браузере. В отличие от
\texttt{sessionStorage}, данные не стираются при закрытии вкладки. Объём ---
около 5~МБ на источник (origin). Хранятся только строки, поэтому объекты
сериализуются через \texttt{JSON.stringify} и читаются через
\texttt{JSON.parse}.

\begin{lstlisting}[caption={Сохранение и чтение данных},captionpos=b]
localStorage.setItem('resume', JSON.stringify(data));
const data = JSON.parse(localStorage.getItem('resume'));
\end{lstlisting}

Доступ к \texttt{localStorage} может быть запрещён (приватный режим,
настройки браузера) --- вызовы оборачиваются в \texttt{try/catch}, при сбое
приложение продолжает работу без сохранения. Все данные \texttt{localStorage}
не покидают браузер пользователя --- это естественный способ хранить личную
информацию, не передавая её на сервер.

\section{Чтение файлов: \texttt{FileReader} и \texttt{canvas}}

Элемент \texttt{<input type="file">} даёт доступ к выбранному пользователем
файлу. Объект \textbf{FileReader} читает файл асинхронно; метод
\texttt{readAsDataURL} возвращает содержимое в виде \textit{data-URL} ---
строки \texttt{data:image/...;base64,...}, которую можно подставить в
атрибут \texttt{src} изображения.

Чтобы фотография не занимала лишнее место в \texttt{localStorage}, она
уменьшается через элемент \texttt{<canvas>}: изображение рисуется на холсте
меньшего размера методом \texttt{drawImage}, а \texttt{canvas.toDataURL}
отдаёт сжатый JPEG.

\section{Печать страницы: \texttt{window.print()} и \texttt{@media print}}

Метод \texttt{window.print()} открывает системный диалог печати. Современные
браузеры предлагают в нём виртуальный принтер <<Сохранить как PDF>>, что
позволяет экспортировать страницу в PDF \textbf{без единой сторонней
библиотеки}. Имя сохраняемого файла берётся из \texttt{document.title}.

Внешний вид печатной версии задаётся медиавыражением \texttt{@media print}:
в нём скрываются интерфейсные элементы (панель управления, подсказки),
убираются тени и фоновые подложки, а правило \texttt{break-inside:~avoid}
запрещает разрывать смысловые блоки между страницами. Так из одной и той же
разметки получаются и экранная, и печатная версии документа.

\begin{lstlisting}[caption={Правила только для печати},captionpos=b]
@media print {
    .toolbar { display: none; }
    .entry { break-inside: avoid; }
}
\end{lstlisting}

\section{CSS-анимации}

\textbf{CSS-анимация} описывается правилом \texttt{@keyframes}, задающим
ключевые кадры (значения свойств в процентах времени), и свойством
\texttt{animation}, которое привязывает набор кадров к элементу и задаёт
длительность, функцию плавности и число повторов. Анимация выполняется
браузером и не требует JavaScript --- скрипт лишь добавляет или снимает
CSS-класс, запускающий нужный эффект.

\begin{lstlisting}[caption={Описание и применение анимации},captionpos=b]
@keyframes saved-flash {
    0%   { background-color: rgba(34, 197, 94, 0.45); }
    100% { background-color: transparent; }
}
.is-saved { animation: saved-flash 420ms ease-out; }
\end{lstlisting}

Медиавыражение \texttt{@media (prefers-reduced-motion: reduce)} позволяет
отключить анимации для пользователей, чувствительных к движению.

\section{Эффект Material Wave (ripple)}

\textbf{Material Wave}, или \textit{ripple} (<<волна>>), --- характерный
приём дизайн-системы Google \textbf{Material Design}: при клике из точки
касания расходится круговая волна, дающая визуальный отклик на действие.
Эффект реализуется без библиотек: по событию \texttt{pointerdown} скрипт
вычисляет координаты клика относительно элемента, создаёт \texttt{<span>}
нужного диаметра и помещает его в точку касания; CSS-анимация увеличивает
круг и плавно уводит его в прозрачность, после чего \texttt{<span>}
удаляется. Элемент-хост получает свойства \texttt{position:~relative} и
\texttt{overflow:~hidden}, чтобы волна не выходила за его границы.

\endinput
